\section{Related work}
\citet{bouthillier2021} derive variance inflation from correlated repeated trainings of one pipeline on one task and measure each source's variance in five case studies, whereas we apply covariance accounting across benchmarks within a replicate run. \citet{reimers2017}, \citet{zhou2020}, and \citet{madaan2024} study seed-related evaluation variability, \citet{dodge2020} examine whether favourable initialisations transfer across tasks, and \citet{jordan2024} separates example-specific training variance from distribution-level variation and compares accuracies across related evaluation sets. Among the studies we located, we find none that estimates the covariance of run noise between benchmarks or an effective benchmark count from that covariance.

Our cross-half construction follows repeated-measurement and variance-component ideas from \citet{spearman1904}, \citet{cronbach1951}, and \citet{searle1992}. The quantitative-genetics account of covariance structure \citep{martin1977} is an analogy that our model doesn't carry as far as separate variance components. \citet{cheverud1988} motivates our phenotypic plug-in comparison (Table~\ref{tab:prediction}), and that comparison fails here on the transfer of correlation structure.

Our few-cluster inference follows the concerns studied by \citet{cameron2008} and \citet{mackinnon2017}. \citet{heineman2025} measure per-benchmark signal and noise on DataDecide and OLMo checkpoints and relate the ratio to decision accuracy, but don't estimate covariance between benchmarks. \citet{cheverud2001}, \citet{nyholt2004}, and \citet{li2005} count effective independent tests to adjust thresholds, whereas we estimate the run covariance of a fixed battery (Appendix~\ref{app:design}).

\section{Discussion and limitations}
We read the accuracy result as a bound rather than as evidence of independence. The full-battery accuracy interval includes one at a design that would detect a true ratio of 1.10 in 0.135 of replicates, so the primary analysis gives a bound of 1.157. The exploratory cuts whose intervals exclude one don't settle it either way, without BoolQ at 1.558 with interval \ci{1.412}{1.687}, at the adjacent checkpoint, and on the held-out tasks. On margins the covariance sits almost entirely within scoring formats.

Apart from the registered test, which failed at its registered scope, and the two PolyPythias rules fixed before scoring, our empirical study remains exploratory. We didn't score the 375 DataDecide runs on the second bank, which leaves the 1.244 untested on fresh items and the shared-item check unrun in the family that supplies our headline. The passage-aware split moves both estimates within re-split variation, and that null follows from BoolQ's passage structure rather than from a sensitive test. A scalar item effect shared across benchmarks and repeated across halves would enter the numerator whole, and shares of 0.124 and 0.009 of item-noise variance reproduce 1.244 and 1.078, an item effect we haven't measured and can't separate from seed covariance with the diagnostics we ran (Appendix~\ref{app:calibration}).\outcome{\Rtwocase}{1}{R2 not supported}{ The cross-bank check in PolyPythias constrains that item effect in one family, and it can't reach an item effect that the two banks share.} A scalar component shared by every benchmark would also enter the cross-format pairs, and those average 0.007 with a diagnostic of 0.921, a check that constrains a uniform component and leaves one that follows scoring format untouched. At a share of 0.05 only 0.007 of simulated replicates reach the observed margin value. We make no claim about that mechanism, since repaired margin proxies fall below one even in simulations without a separate competence variable, apart from the clipped-drop weighting at 0.889 against a simulated fifth percentile of 0.908.

We cluster on recipes, but the release shares one seed label across every recipe inside a size band, and a run effect following the band would break that choice and take simulated margin coverage to 0.821 at a quarter of latent variance and 0.580 at a half. However, a permutation test that aligns runs by batch slot bounds that effect at 0.044 of the seed covariance on margins, where simulated coverage is 0.941 at a share of 0.05, and that test's power is only 0.490 at the bound. Within a single size we find the evidence weaker, as the cluster-$t$ 1B margin interval of 0.811 to 1.628 includes one (Appendix~\ref{app:supplementary}).

For practice we recommend paired replicate scores across the chosen battery, which give the aggregate seed standard deviation directly, and a benchmark-removal sensitivity check. Out of sample on margins, independence predicts a configuration's measured ratio with a mean absolute log error of 0.339 against 0.536 for the plug-in and 0.512 for our decomposition (Appendix~\ref{app:transport}). Our estimated ratio is therefore a measurement of this battery and helps only where replicates are missing and the target battery's own ratio is close. A scaled independence standard error, with a ratio estimated on the same margin population, returns the simulated false-positive rate from 0.113 to 0.051, and carried to a battery whose own ratio is 1.5, that scaling leaves the error at 0.101. Under an equicorrelated seed covariance at 125 configurations, three runs per configuration let the margin interval exclude one in 0.997 of replicates, while accuracy needs eight runs to reach 0.803.

\label{maintext:end}
\section*{AI use statement}
We used generative AI tools to propose and refine hypotheses, critique the estimator and simulation design, assist with analysis code, revise manuscript text, survey literature, and contribute to the technical and editorial review of this version. The study includes synthetic calibration and sensitivity simulations, whose role differs from the released model outputs used for the empirical analysis. The authors remain responsible for verifying AI-assisted claims, code, citations, and the final submission, including an accurate disclosure of AI involvement in generating simulation data.

\section*{Ethics statement}
This study analyses released model outputs and introduces no new participant recruitment or data collection. Its intended use is to improve uncertainty reporting in model evaluation. An independence approximation used outside the studied battery could understate uncertainty, which keeps the recommendations conditional on the measured population and score definition. Use of derived artifacts must retain the applicable source attribution and licensing requirements.

\section*{Reproducibility statement}
Section~\ref{sec:estimator} and Appendix~\ref{app:derivations} define the estimator and its assumptions. Appendix~\ref{app:implementation} specifies the interval calculation and data reduction. Appendix~\ref{app:design} records the populations and planned checks. Appendix~\ref{app:supplementary} reports later sensitivity analyses alongside the primary results. The released archive also holds an extended supplement with the full-length versions of the appendix sections that this version compresses. The primary calculation uses 4,999 bootstrap draws, 2,000 calibration replicates, 500 permutation replicates, 50 item re-splits, and master seed 20260101, and reproduction requires the per-item reductions, item ordering and half assignments, selected checkpoint identifiers, and the corresponding analysis scripts. We release all four of them alongside the paper. The original analysis plan lacks an independently corroborated timestamp. Its proposed screening split was never implemented either. The additional simulation, transport, and checkpoint analyses require their own execution records. A separate implementation run on Kaggle reproduces the primary estimates, the BoolQ-removal estimates, and their intervals to three decimals from the released per-item scores, a run that also produced the passage-aware split, the coverage simulation, the operational prediction comparison, and the seed-matched pair analysis. The PolyPythias transfer, five-size rescore, and checkpoint results come from earlier execution records that this revision didn't rerun, and the loss-proxy rescoring wasn't run. The registered test ships with its protocol, the dated amendments that record each change to it, the frozen hash of its request file, and the kernels that scored and analysed it, and the last twenty runs at 1B were scored off Kaggle on three RTX 3090 cards. The PolyPythias battery and the second bank were scored in float32 with transformers 4.57.1 on Kaggle T4 cards, each run checking its cached scores against an uncached reference on 24 items to within 0.001 nats, and the release ships the frozen hashes of both request files with the kernels that scored them. The protocol's commit times aren't corroborated by an independent source. Separate CPU kernels on the released runs produced the recipe-decision, seed-count, battery-size, and item-count analyses.

